\textbf{Determinação do modelo de erro:}

O erro varia \emph{linearmente} com a leitura $R$. Usando os dois pontos dados:
\[
(R_1, E_1)=(5,1),\quad (R_2, E_2)=(80,4).
\]

A equação da reta:
\[
m=\frac{E_2-E_1}{R_2-R_1}=\frac{4-1}{80-5}=\frac{3}{75}=0{,}04.
\]
\[
E(R)=E_1+m(R-R_1)=1+0{,}04(R-5)=0{,}04R+0{,}8\ \mathrm{V}.
\]

Verificação: $E(5)=0{,}04\times5+0{,}8=1{,}0$ $\checkmark$; $E(80)=0{,}04\times80+0{,}8=4{,}0$ $\checkmark$.

\textbf{(a) Tensão correta $T(R)=R-E(R)$:}
\[
T(R)=R-(0{,}04R+0{,}8)=0{,}96R-0{,}8.
\]

\begin{center}
\begin{tabular}{llll}
\hline
Leitura $R$ (V) & Erro $E(R)$ (V) & Tensão correta $T$ (V) & Erro $\%$ \\
\hline
$80$ & $4{,}00$ & $76{,}00$ & $5{,}26\%$ \\
$100$ & $4{,}80$ & $95{,}20$ & $5{,}04\%$ \\
$50$ & $2{,}80$ & $47{,}20$ & $5{,}93\%$ \\
$1$ & $0{,}84$ & $0{,}16$ & $525\%$ \\
$35$ & $2{,}20$ & $32{,}80$ & $6{,}71\%$ \\
$10$ & $1{,}20$ & $8{,}80$ & $13{,}64\%$ \\
\hline
\end{tabular}
\end{center}

\textbf{(b) Observações:}
\begin{itemize}
\item O erro percentual é aproximadamente constante ($\approx5\%$) para leituras altas.
\item Para leituras baixas (ex: $R=1\,\mathrm{V}$), o erro percentual é enorme (525\%) devido ao erro fixo de offset ($0{,}8$).
\item O voltímetro é confiável apenas para leituras acima de $\approx20\,\mathrm{V}$.
\end{itemize}